
%\needspace{4\baselineskip}
\color{uniblue}{\subsection{Блокировка при неисправности цепей напряжения (БНН)}}
\color{black}

\begin{enumerate}[label=\arabic{section}.\arabic{subsection}.\arabic*, labelsep=4pt, leftmargin=0pt, itemindent=57pt]

\item
Блокировка при неисправности цепей напряжения используется для блокировки органов защит, которые могут работать неправильно при нарушениях в цепях напряжения.

\item
Функция <<БНН>> включает в себя четыре различных алгоритма:
\begin{itemize}
    \item контроль напряжения небаланса U БНН (dU0), полученного из определенной комбинации мгновенных значений фазных напряжений <<звезды>> и напряжения $3U_0$ <<разомкнутого треугольника>>;
    \item контроль снижения линейных напряжений и повышения напряжения ОП (U2minU1p) в режиме работы секции;
    \item контроль включенного состояния защитных автоматических выключателей во вторичных цепях напряжения при вкаченном ВЭ ТН;
    \item контроль третьей гармоники напряжения нулевой последовательности (3U03f).
\end{itemize}

\item
Срабатывание функции <<БНН>> с режимом работы контроля небаланса (уставка SGF2) происходит при возникновении неисправности во вторичных цепях ТН, при этом баланс напряжений обмоток <<звезды>> и <<треугольника>> нарушается. Алгоритм рассчитывает среднее значение суммы фазных напряжений за полпериода и сравнивает его со средним измеренным значением $3U_0$ для того же полупериода. Если разность двух параметров превышает заданную уставку, функция срабатывает и активируется соответствующий выходной сигнал <<БНН:~Пуск dU0>>.
\item
Напряжение небаланса рассчитывается по формуле:
% ФОРМУЛА
\begin{equation*}
    dU_0=\lvert 3U_{0\textsl{вычисл.}} - \frac{3U_{0\textsl{вычисл.ном.}}}{3U_{0\textsl{изм.ном.}}}3U_{0\textsl{изм.}}  \rvert,
\end{equation*}
\begin{eqexpl}[25mm]
    \item{$3U_{0\textsl{вычисл.}}$} сумма фазных напряжений <<звезды>>;
    \item{$3U_{0\textsl{изм.}}$} измеренное напряжение $3U_0$ <<разомкнутого треугольника>>;
    \item{$3U_{0\textsl{вычисл.ном.}}$} номинальное напряжение основной обмотки ТН;
    \item{$3U_{0\textsl{изм.ном.}}$} номинальное напряжение $3U_0$ обмотки <<разомкнутого треугольника>> ТН.
\end{eqexpl}
\item
Симметричность вторичных напряжений, подводимых к устройству, определяется уровнем напряжения обратной последовательности $U_2$ и наличием линейного напряжения. Срабатывание функции БНН с режимом контроля от $U_2$ будет формироваться при увеличении напряжения ОП выше уставки с фиксированной выдержкой времени Т3 составляющей 20 мс, или снижении одного из трех линейных напряжений ниже уставки при включенном ВВ или СВ по истечении выдержки времени таймера Т2.

\item
Для алгоритма контроля исправности вторичных цепей ТН по включенному положению защитных выключателей в цепях ТН используются блок-контакты защитных выключателей (<<звезда>>/<<треугольник>>) при вкаченном положении ВЭ ТН. Сигнал с блок-контактов АВ ТН заводится на дискретные входы устройства <<Положение SF1>> и <<Положение SF2>>. При отключенном положении АВ ТН или выкаченного положения ВЭ ТН формируется выходной сигнал <<БНН: АВ ТН откл>>. 

\item
Уставкой SGF3 осуществляется ввод в работу алгоритма контроля утроенного напряжения НП третьей гармоники, который отслеживает обрыв в цепях <<разомкнутого треугольника>>. Срабатывание ИО информируется сигналом ИО «БНН:3U03f<», реагирующим на снижение уровня напряжения третьей гармоники в цепях <<разомкнутого треугольника>> ниже уставки. В данном случае учитывается тот факт, что даже в обычном режиме работы из-за нелинейности ТН в схемах <<разомкнутого треугольника>> наблюдается заметный уровень третьей гармоники, составляющий около 0,5 В. При обрыве, уровень третьей гармоники значительно снижается, теоретически приближаясь к нулю. Для исключения ложного срабатывания в нормальных и аварийных режимах, не связанных с неисправностью в цепях <<разомкнутого треугольника>>, выполнен контроль уровня минимального напряжения и контроль напряжения ОП. При выявлении неисправностей в цепях <<разомкнутого треугольника>> формируется сигнал неисправности в цепях напряжения.

\item
Наличие включенного ВВ или СВ, при введенной уставке SGF4 указывает о том, что секция введена в работу (активируется выходной сигнал <<БНН: Секция в работе>>).

\item
Формирование сигнала сигнализации о неисправности ЦН <<БНН: Неисправность ЦН>> выполняется по истечении выдержки времени таймера T1 при выполнении одного из следующих условий:
\begin{itemize}
    \item наличие условий формирования сигнала срабатывания БНН <<БНН: Срабатывание>>;
    \item наличие утроенного напряжения НП третьей гармоники.
\end{itemize}

\item
Формирование сигнала блокировки ЗМН <<БНН: Блокировка ЗМН> выполняется когда:
\begin{itemize}
    \item ВВ и СВ отключен;
    \item напряжения ОП выше уставки в течении 20 мс;
    \item присутствует неисправность в цепях ТН;
    \item отключен один из автоматических выключателей <<звезда>>/<<треугольник>> или ВЭ ТН выкачен.
\end{itemize}

\item
Ввод функции <<БНН>> в работу на этапе параметрирования устройства осуществляется программным переключателем SGF1 <<Ввод\_функции>>. Информация о введенной ступени отображается выходным сигналом <<БНН: Ввод>> на выходе функции.
\item
Функция <<БНН>> может быть выведена из работы оперативно путем активации сигнала <<ОВ БНН>>, что характеризуется активным состоянием сигнала <<БНН: Оперативный вывод>>.

\item
Алгоритм работы ступени <<БНН>> выполнен в соответствии с рисунком \ref{bnn:img1}. В таблице \ref{bnn:tbl1} приведены параметры, необходимые для настройки функции <<БНН>>.

\small
\begin{longtable}{|>{\centering\arraybackslash}m{5.3cm}|>{\centering\arraybackslash}m{3.3cm}|>{\centering\arraybackslash}m{4.2cm}|>{\centering\arraybackslash}m{1.8cm}|>{\centering\arraybackslash}m{1cm}|}
\caption{Параметры для настройки функции <<БНН>>\hfill\vspace{-0.5\baselineskip}}\label{bnn:tbl1}\\ 
\hline
\rowcolor{gray!30}
Параметр (Параметр на ИЧМ) & Условное обозначение на схеме & Значение/ Диапазон & Единица измерения & Шаг \\ 
\hline
\endfirsthead
\caption*{\hspace{3pt}\emph{Продолжение таблицы \ref{bnn:tbl1}\hfill\vspace{-0.5\baselineskip}}} \\ % сделано по ГОСТ 2.105 п.6.8.7
\hline
\rowcolor{gray!30}
Параметр (Параметр на ИЧМ) & Условное обозначение на схеме & Значение/ Диапазон & Единица измерения & Шаг \\ 
%\hline
\endhead
%\hline
\endfoot
%\hline
\endlastfoot
%==+t1*RVTR1|VT_LVVTVTR
\centering Ввод функции в работу (Ввод\_функции) & \centering SGF1 & \centering 0 = Не предусмотрено\\1 = Предусмотрено & \centering -- & \centering \arraybackslash -- \\
\hline
\centering Уровень небаланса по напряжению (dU0ср) & \centering Uф--3U0> & \centering 0,0 ... 150,0 & \centering В & \centering \arraybackslash 0,1 \\
\hline
\centering Напряжение срабатывания ОП (U2ср) & \centering U2> & \centering 2,0 ... 60,0 & \centering В & \centering \arraybackslash 0,1 \\
\hline
\centering Напряжение срабатывания НП третьей гармоники (3U0h3ср) & \centering 3U03f< & \centering 0,05 ... 10,00 & \centering В & \centering \arraybackslash 0,01 \\
\hline
\centering Напряжение срабатывания (Uср) & \centering U< & \centering 2,0 ... 100,0 & \centering В & \centering \arraybackslash 0,1 \\
\hline
\centering Ввод в работу контроля небаланса (Контроль\_dU0) & \centering SGF2 & \centering 0 = Не предусмотрено\\1 = Предусмотрено & \centering -- & \centering \arraybackslash -- \\
\hline
\centering Ввод в работу контроля 3 гармоники напряжения НП (Контроль\_3гарм.3U0) & \centering SGF3 & \centering 0 = Не предусмотрено\\1 = Предусмотрено & \centering -- & \centering \arraybackslash -- \\
\hline
\centering Ввод контроля включенного состояния секции (Контроль\_секции) & \centering SGF4 & \centering 0 = Не предусмотрено\\1 = Предусмотрено & \centering -- & \centering \arraybackslash -- \\
\hline
\centering Выдержка времени срабатывания (Tср) & \centering T1 & \centering 0 ... 100000 & \centering мс & \centering \arraybackslash 10 \\
\hline
\centering Выдержка времени срабатывания сигнализации (Tсигн) & \centering T2 & \centering 0 ... 100000 & \centering мс & \centering \arraybackslash 10 \\
\hline
%===t1
\end{longtable}
\normalsize

\vspace{3mm}
\begin{figure}[H]
\centering
\includegraphics[width=1\textwidth,height=1\textheight,keepaspectratio]{\fbpath/204. БНН ТН/_latex/img/bnn.pdf}
\captionof{figure}{Функциональная схема <<БНН>>}\label{bnn:img1}
\end{figure}

\end{enumerate}
